\documentclass{osucourse}
\course{1136}

%% Not in the university's course data, so these come
%% from the published sheet this file was imported from.
\SetCourseField{title}{Measurement \& Geometry for Teachers}
\SetCourseField{credits}{5 credits}
\SetCourseField{description}{This course is the second in a two semester sequence for teachers of elementary and middle grade students. This course focuses on concepts of measurement and geometry, including modern and historical perspectives.}
\SetCourseField{prereq}{A grade of C$-$ or above in “Number and Operations for Teachers” (Math 1135)}

\begin{document}

\CatalogDescription
\PrerequisitesAndExclusions

\begin{Purpose}
The course consists of fundamental topics in Euclidean geometry starting with measurement. This includes the concepts of length, area, volume, angles, units of measurement, precision and error.

The basic properties of two and three dimensional geometric shapes and their relationships are a central part of the course. Special emphasis is put on geometric reasoning through problem solving, including unknown angle, length, area, and volume. The course also covers topics on transformations in the plane, symmetries, congruence, and similarity. Some geometric constructions and basic geometric proofs are included.

Additional topics include an introduction to functions and equations, primarily in the linear case, and a brief introduction to probability.
\end{Purpose}

\begin{Text}
\emph{Mathematics for Elementary Teachers, with Activity Manual}, 4\textsuperscript{th} Edition, by Sybilla Beckmann, Pearson, ISBN for the package is 9780321836715 (loose-leaf).
\end{Text}

\begin{SupplementalTexts}
• \emph{Geometric Structures: An Inquiry-Based Approach for Prospective Elementary and Middle School Teachers}, by Douglas Aichele and John Wolfe, Pearson, ISBN 9780131483927 • \emph{Elementary Geometry for Teachers}, by Thomas Parker and Scott Baldridge, Sefton-Ash Publishing, ISBN 9780974814056
\end{SupplementalTexts}

\begin{TopicsList}
\item Measurement
\item Planar shapes
\item Polyhedra
\item Plane geometry
\item Transformations in the plane, congruence, symmetry
\item Linear equations and graphs
\item Algebra and linear equations
\item Probability
\end{TopicsList}

\end{document}
